\documentclass[11pt,a4paper]{article}

\usepackage[margin=2.6cm]{geometry}
\usepackage[T1]{fontenc}
\usepackage{lmodern}
\usepackage{microtype}
\usepackage{amsmath,amssymb}
\usepackage{booktabs}
\usepackage{enumitem}

\setlength{\parindent}{0pt}
\setlength{\parskip}{6pt}
\setlist{nosep,leftmargin=1.6em}

\title{Outcome Imputation in \texttt{predict.crc\_fit()}}
\author{Internal mathematical note}
\date{28 August 2026}

\begin{document}

\maketitle

\section{Purpose and scope}

This note describes the mathematical calculations used by
\texttt{predict.crc\_fit()} to estimate the total value of the modeled outcome.
The method combines conditional imputation for observed units with prediction
for units belonging to the all-zero capture cells. It applies to exact,
interval-censored, right-censored, and missing positive count outcomes and to
both outcome distributions supported by the package.

The calculation is performed after the EM algorithm has converged. All model
parameters and posterior state weights below are therefore evaluated at their
final fitted values. Hats are omitted where doing so does not cause ambiguity.

\section{State and outcome notation}

Let $i\in S$ index an observed unit and let $s\in\mathcal S_i$ index a
candidate state of that unit. A state records the candidate true values of
misclassified categorical variables and the latent class. Write
\[
  \tau_{is}
  =\Pr(s\mid\text{observed information for unit }i)
\]
for its final posterior weight. These weights satisfy
\[
  \tau_{is}\geq0,
  \qquad
  \sum_{s\in\mathcal S_i}\tau_{is}=1.
\]

For state $s$, the outcome regression provides the linear predictor and mean
\[
  \eta_{is}=z_{is}^{\mathsf T}\beta,
  \qquad
  \mu_{is}=\exp(\eta_{is}).
\]
The parameter $\mu_{is}$ is the mean of an underlying untruncated count $W$.
The modeled positive outcome is
\[
  V=W\mid W\geq1.
\]
The underlying distribution is either
\[
  W\sim\operatorname{Poisson}(\mu)
\]
or
\[
  W\sim\operatorname{NB}(\mu,r),
  \qquad r=\frac{1}{\sigma},
\]
under the mean--size parameterization
\[
  \operatorname{E}(W)=\mu,
  \qquad
  \operatorname{Var}(W)=\mu+\sigma\mu^2.
\]

Define the survival function
\[
  S(q;\mu)=\Pr(W>q\mid\mu).
\]
The implementation evaluates $\log S(q;\mu)$ directly through the upper-tail
interfaces of the Poisson and negative-binomial distribution functions.

\section{Shifted-distribution identities}

Conditional means over censored ranges can be evaluated without explicitly
summing $w\Pr(W=w)$. For a Poisson random variable,
\[
  w\Pr(W=w;\mu)
  =\mu\Pr(W^\star=w-1;\mu),
\]
where $W^\star\sim\operatorname{Poisson}(\mu)$. Consequently,
\[
  \sum_{w=l}^{u}w\Pr(W=w)
  =\mu\Pr(l-1\leq W^\star\leq u-1).
\]

For a negative-binomial random variable with size $r$ and mean $\mu$,
\[
  w\Pr(W=w;\mu,r)
  =\mu\Pr(W^\star=w-1;\mu^\star,r^\star),
\]
where
\[
  r^\star=r+1,
  \qquad
  \mu^\star=\mu\frac{r+1}{r}.
\]
This choice preserves the probability parameter of the negative-binomial
distribution. The implementation calls this the \emph{shifted} distribution.
Let $S^\star(q;\mu)$ denote its survival function. For the Poisson model,
$S^\star$ has the same parameters as $S$; for the negative-binomial model, it
uses $(\mu^\star,r^\star)$.

The preceding identities imply
\begin{align*}
  \operatorname{E}\{W\mathbf 1(W\geq l)\}
    &=\mu S^\star(l-2;\mu),\\
  \operatorname{E}\{W\mathbf 1(l\leq W\leq u)\}
    &=\mu\{S^\star(l-2;\mu)-S^\star(u-1;\mu)\}.
\end{align*}

\section{Conditional imputation for observed units}

For every observed unit and candidate state, the function constructs
\[
  v_{is}^{\mathrm{imp}}
  =\operatorname{E}(V_i\mid\mathcal O_i,s),
\]
where $\mathcal O_i$ denotes the exact, censored, or missing outcome
information available for unit $i$.

\subsection{Missing outcome}

When the outcome is missing, the only outcome information is the positive
support of $V$. Since the zero value contributes nothing to
$\operatorname{E}(W)$,
\[
  \operatorname{E}(V\mid\mu)
  =\operatorname{E}(W\mid W\geq1)
  =\frac{\mu}{\Pr(W\geq1)}
  =\frac{\mu}{S(0;\mu)}.
\]
The implementation evaluates this as
\[
  \exp\{\log\mu-\log S(0;\mu)\},
\]
which avoids overflow caused by separately calculating the reciprocal of a
very small survival probability.

\subsection{Exact outcome}

For an exact observed value $v_i$,
\[
  v_{is}^{\mathrm{imp}}
  =\operatorname{E}(V_i\mid V_i=v_i,s)
  =v_i.
\]
This value is independent of the candidate state. Consequently, posterior
averaging preserves the observation because
\[
  \sum_{s\in\mathcal S_i}\tau_{is}v_i=v_i.
\]

\subsection{Right-censored outcome}

Suppose the available information is $V_i\geq l_i$. Because $l_i\geq1$,
conditioning additionally on zero truncation is redundant. The conditional
mean is
\begin{align*}
  \operatorname{E}(W\mid W\geq l)
  &=\frac{\operatorname{E}\{W\mathbf 1(W\geq l)\}}
          {\Pr(W\geq l)}\\
  &=\mu\frac{S^\star(l-2;\mu)}{S(l-1;\mu)}.
\end{align*}
The log-scale expression used in the function is
\[
  \exp\{\log\mu+\log S^\star(l-2;\mu)-\log S(l-1;\mu)\}.
\]
When $l=1$, $S^\star(-1;\mu)=1$, and the result reduces to the ordinary
zero-truncated mean.

\subsection{Interval-censored outcome}

Suppose that $l_i\leq V_i\leq u_i$, with both endpoints included. The
denominator of the conditional expectation is
\[
  \Pr(l\leq W\leq u)
  =S(l-1;\mu)-S(u;\mu).
\]
Using the shifted-distribution identity, its numerator is
\[
  \operatorname{E}\{W\mathbf 1(l\leq W\leq u)\}
  =\mu\{S^\star(l-2;\mu)-S^\star(u-1;\mu)\}.
\]
It follows that
\[
  \operatorname{E}(W\mid l\leq W\leq u)
  =\mu
   \frac{S^\star(l-2;\mu)-S^\star(u-1;\mu)}
        {S(l-1;\mu)-S(u;\mu)}.
\]

Both probability differences are evaluated on the log scale. If
$a=\log A$ and $b=\log B$ with $A\geq B$, then
\[
  \log(A-B)=a+\log\{-\operatorname{expm1}(b-a)\}.
\]
This is the operation implemented by \texttt{log\_sub\_crc()}. It avoids
catastrophic cancellation when two right-tail probabilities are close.

\section{Observed-unit contribution}

Each expanded state represents a fractional expected unit with population
contribution $\tau_{is}$. Its estimated outcome contribution is
\[
  \tau_{is}v_{is}^{\mathrm{imp}}.
\]
Summing over states and observed units gives
\[
  \widehat T_{\mathrm{obs}}
  =\sum_{i\in S}\sum_{s\in\mathcal S_i}
    \tau_{is}v_{is}^{\mathrm{imp}}.
\]
The posterior weights already incorporate the capture model,
misclassification probabilities, and outcome likelihood. Thus, imputation for
a censored or missing outcome is averaged over uncertainty about its true
categorical values and latent class.

\section{Contribution from completely unobserved units}

Let $c$ index a complete capture cell with the all-zero capture history. The
fitted capture model supplies its expected population count
\[
  m_c=\exp(x_c^{\mathsf T}\theta).
\]
All categorical variables required by the outcome model are present in the
capture-cell grid. Therefore, the outcome design row $z_c$ and untruncated
outcome mean
\[
  \mu_c=\exp(z_c^{\mathsf T}\beta)
\]
can be constructed for every all-zero cell. Since no outcome is observed for
these units, their expected positive outcome is
\[
  \operatorname{E}(V\mid c)
  =\frac{\mu_c}{S(0;\mu_c)}.
\]
The estimated outcome total contributed by cell $c$ is consequently
\[
  m_c\frac{\mu_c}{S(0;\mu_c)}.
\]
The implementation rejects non-finite or non-positive values of $\mu_c$ and
evaluates the zero-truncated mean on the log scale.

\section{Overall and subpopulation estimators}

Combining observed and completely unobserved contributions gives
\[
  \widehat T
  =\widehat T_{\mathrm{obs}}
   +\sum_{c:\,\boldsymbol y_c=\boldsymbol 0}
      m_c\frac{\mu_c}{S(0;\mu_c)}.
\]
The corresponding population-size estimator is
\[
  \widehat N
  =\sum_{i\in S}\sum_{s\in\mathcal S_i}\tau_{is}
   +\sum_{c:\,\boldsymbol y_c=\boldsymbol 0}m_c
  =|S|+\sum_{c:\,\boldsymbol y_c=\boldsymbol 0}m_c.
\]

For a subpopulation $A$ defined by categorical variables, the latent class,
or both, the function restricts the preceding sums to states and all-zero
cells belonging to $A$:
\begin{align*}
  \widehat N_A
  &=\sum_{i\in S}\sum_{s\in\mathcal S_i}
      \tau_{is}\mathbf 1(s\in A)
    +\sum_{c:\,\boldsymbol y_c=\boldsymbol 0}
      m_c\mathbf 1(c\in A),\\
  \widehat T_A
  &=\sum_{i\in S}\sum_{s\in\mathcal S_i}
      \tau_{is}v_{is}^{\mathrm{imp}}\mathbf 1(s\in A)
    +\sum_{c:\,\boldsymbol y_c=\boldsymbol 0}
      m_c\frac{\mu_c}{S(0;\mu_c)}\mathbf 1(c\in A).
\end{align*}
These subgroup estimates add to the corresponding overall estimates when the
grouping variables define a partition of the modeled population.

\section{Summary of the implemented cases}

\begin{center}
\begin{tabular}{@{}ll@{}}
\toprule
Outcome information & Imputed state value\\
\midrule
Missing
  & $\mu/S(0;\mu)$\\
Exact $v$
  & $v$\\
Right-censored at $l$
  & $\mu S^\star(l-2;\mu)/S(l-1;\mu)$\\
Interval $[l,u]$
  & $\mu\{S^\star(l-2;\mu)-S^\star(u-1;\mu)\}/
     \{S(l-1;\mu)-S(u;\mu)\}$\\
\bottomrule
\end{tabular}
\end{center}

\end{document}
